% III. Проектиране на системата.
% Структура, интерфейси, формати на данните, разположение в паметта.
\section{Проектиране на системата}
\label{sec:design}

Системата е разделена на четири слоя. Най-долният съдържа представянето на данните:
облак от точки и растерно изображение. Над него стои пространственият индекс с единен
интерфейс за заявка по $k$ най-близки съседи и за заявка по радиус. Третият слой
съдържа стъпките на обработката, всяка от които зависи само от долните два. Най-горният
слой са приложенията: демонстрацията, измерванията и проверките. Разделянето е показано
на Фиг.~\ref{fig:layers}.

\begin{figure}[H]
\centering
\begin{tikzpicture}[node distance=0.55cm and 0.8cm]
  \node[block, minimum width=13cm, fill=gray!8] (apps)
    {\textbf{Слой 4: приложения}\\[2pt]
     \code{apps/demo.cpp} \quad \code{benchmarks/bench.cpp} \quad \code{tests/}};

  \node[block, minimum width=13cm, below=of apps] (stages)
    {\textbf{Слой 3: стъпки на обработката}\\[2pt]
     прореждане \quad съвместяване \quad изглаждане \quad градиент \quad
     потискане на немаксимуми \quad сегментация};

  \node[accent, minimum width=13cm, below=of stages] (index)
    {\textbf{Слой 2: пространствен индекс}\\[2pt]
     \code{KdTree}: \code{knn}, \code{radius\_search}, \code{nearest}, \code{memory\_bytes}};

  \node[block, minimum width=13cm, below=of index, fill=blue!8] (data)
    {\textbf{Слой 1: представяне на данните}\\[2pt]
     \code{PointCloud} \quad \code{Image8} \quad \code{ImageF} \quad
     \code{RigidTransform}, \code{Mat3}};

  \draw[line] (apps) -- (stages);
  \draw[line] (stages) -- (index);
  \draw[line] (index) -- (data);
  \draw[line, dashed] (stages.west) to[out=180,in=180] node[left, font=\scriptsize]
       {растерните стъпки не ползват индекса} (data.west);
\end{tikzpicture}
\caption{Слоеве на системата. Зависимостите сочат само надолу.}
\label{fig:layers}
\end{figure}

\subsection{Разположение на данните в паметта}

Разположението на данните е проектно решение, не подробност. Облакът се съхранява като
три отделни непрекъснати масива, по един за всяка координата, а не като един масив от
тройки. Класът \code{PointCloud} пази \code{std::vector<float>} за $x$, за $y$ и за $z$
и дава пряк достъп до трите указателя, защото векторизираният път ги ползва направо.

Причината е достъпът при заявка. Вътрешният цикъл на обхождането на лист чете една
координата на много точки последователно. При разположение по компонента това е
последователен достъп с крачка 4 байта, тоест едно четене на кеш-линия обслужва
шестнадесет точки. При масив от тройки крачката е 12 байта и същият цикъл докосва три
пъти повече кеш-линии, а векторизацията изисква разместване на регистрите. Цената на
избора е по-неудобен достъп до отделна точка: методът \code{point(i)} събира тройката от
три различни масива. Затова той е на границата на интерфейса, а не във вътрешните
цикли.

\subsection{Интерфейс на индекса}

Интерфейсът на индекса е умишлено тесен: построяване от облак, заявка за $k$ най-близки
съседи, заявка по радиус, заявка за единствен най-близък съсед и отчет за паметта.
Всичко останало в системата е потребител на този интерфейс, а не негова част.

Двата пътя за обхождане на лист, \term{скаларен} и \term{пакетен}, не са две различни
реализации на индекса, а параметър на една и съща заявка. Решението е методическо. Ако
бяха два класа или два варианта на компилация, сравнението между тях щеше да сравнява и
всичко останало, което ги отличава: различни дървета, различно подреждане на паметта,
евентуално различни версии на кода. Като параметър на метода те работят върху едно и
също дърво, върху една и съща подредба, в едно и също изпълнение на програмата, и
единствената разлика между двете измервания е тялото на цикъла. Това е условието
сравнението да е измерване, а не твърдение.

Заявката за единствен най-близък съсед е отделен метод, а не заявка с $k=1$. Причината е
потребителят ѝ: съвместяването вика точно нея веднъж на точка на всяка итерация, тоест
милиони пъти, и купчината с кандидати, която общият път поддържа, там е чиста излишна
работа. Двата метода дават един и същ резултат и това се проверява.

\subsection{Договор на стъпката}

Всяка стъпка приема вход и конфигурация и връща нов резултат. Стъпките не променят входа
си на място. Решението е за сметка на паметта, но има три последици, които го оправдават:
повтарянето на една стъпка с различни параметри не изисква повторно изпълнение на
предходните, отпада скрита зависимост по ред на изпълнение, която би направила
измерването на отделна стъпка невярно, и стъпката става безопасна за паралелно
изпълнение върху различни входове, без това да е било предвиждано отделно.

Конфигурацията е структура с полета по подразбиране, а не файл. Разгледан беше и
конфигурационен файл, който конвейерът да чете при пускане. Отхвърлен е, защото би
въвел или зависимост от библиотека за разбор на формата, или собствен разбор, който е
код без връзка с предмета на работата. Структурите \code{KdTreeOptions},
\code{IcpOptions}, \code{SceneOptions} и \code{ImageSceneOptions} дават същото, което
файлът би дал, а стойностите по подразбиране са на едно място в заглавния файл, където
се четат заедно с описанието на полето.

\subsection{Поток на данните}

Потокът е показан на Фиг.~\ref{fig:dataflow}. Геометричният и растерният клон са
независими: те споделят слоевете отдолу, но никоя стъпка от единия не е вход на стъпка
от другия. Разделянето е следствие от обхвата: работата изследва цената на заявката за
съседи, а тя се среща само в геометричния клон.

\begin{figure}[H]
\centering
\begin{tikzpicture}[node distance=0.45cm and 0.6cm]
  \node[block] (gen) {генератор на\\изкуствени данни};
  \node[block, left=1.4cm of gen] (io) {PLY / PCD\\PGM / PPM};

  \node[accent, below left=0.9cm and 0.2cm of gen] (cloud) {облак\\от точки};
  \node[accent, below right=0.9cm and 0.2cm of gen] (img) {изображение};

  \node[block, below=of cloud] (voxel) {прореждане с\\вокселна решетка};
  \node[block, below=of voxel] (kd) {построяване на\\k-d дърво};
  \node[block, below=of kd] (icp) {съвместяване\\(ICP)};
  \node[danger, below=of icp] (res3d) {преобразувание,\\отклонение, итерации};

  \node[block, below=of img] (blur) {гаусово\\изглаждане};
  \node[block, below=of blur] (sobel) {градиент\\по Собел};
  \node[block, below=of sobel] (nms) {потискане на\\немаксимуми};
  \node[block, below=of nms] (hyst) {двупрагово\\решение};
  \node[danger, below=of hyst] (edges) {карта на\\контурите};

  \node[block, right=1.2cm of nms] (thr) {праг};
  \node[block, below=of thr] (cc) {свързани\\компоненти};
  \node[danger, below=of cc] (segs) {области с\\описание};

  \draw[line] (io) -- (gen);
  \draw[line] (gen) -- (cloud);
  \draw[line] (gen) -- (img);
  \draw[line] (cloud) -- (voxel);
  \draw[line] (voxel) -- (kd);
  \draw[line] (kd) -- node[right, font=\scriptsize] {заявка за съсед} (icp);
  \draw[line] (icp) -- (res3d);
  \draw[line] (img) -- (blur);
  \draw[line] (blur) -- (sobel);
  \draw[line] (sobel) -- (nms);
  \draw[line] (nms) -- (hyst);
  \draw[line] (hyst) -- (edges);
  \draw[line] (blur.east) to[out=0,in=180] (thr.west);
  \draw[line] (thr) -- (cc);
  \draw[line] (cc) -- (segs);
\end{tikzpicture}
\caption{Поток на данните. Геометричният клон е отляво, растерният отдясно.}
\label{fig:dataflow}
\end{figure}

\subsection{Извън обхвата}

Три решения от първоначалния замисъл са отпаднали и е редно да бъде записано защо, а не
да останат като липса.

Втора реализация на индекса, вокселна решетка като заменяем пространствен индекс, е
изоставена. Вокселната решетка е реализирана, но като \term{стъпка за прореждане}, което
е ролята, в която тя носи полза в този конвейер. Като индекс тя би трябвало да отговаря
на същите заявки като k-d дървото, а при неравномерна плътност, каквато има в изпитваната
сцена, това означава или памет за празни кубчета, или прекалено много точки в едно.
Сравнението между двете щеше да измерва това, което Раздел~\ref{sec:alternatives} вече
установява от устройството им.

Сравнение с еталонна библиотека с отворен код е изоставено. Причината е ограничението за
изграждане: проектът се компилира с компилатор за C++20 и CMake и с нищо друго, а
сравнението би добавило задължителна зависимост заради една таблица. Еталонът за
коректност на индекса е изчерпателното търсене, което е точно по определение и е
реализирано в проекта.

Оценяването на дълбочина по стереодвойка и извличането на описания на околност са
изоставени по обем. Те не участват в измерваната верига: и двете са потребители на
заявката за съседи, а нейната цена вече се измерва пряко.
